\section{The original support action on prime powers, kernel memory and tensor amplification}
\label{sec:tau-prime-memory}

This calculation uses the exact objects and metric of ST1--ST20 in
\emph{What the original observation retains of split support}, result
SZ-20260922-028. It neither substitutes a new source metric nor discards
the original observation kernel. The human mathematical precedents for
the final comparison are Pierre Deligne, \emph{La conjecture de Weil. II},
Publications Math\'ematiques de l'IH\'ES \textbf{52} (1980), 137--252,
\S\S1.4--1.5 and 3.2, and the Grothendieck trace formula he uses in
1.4.5.1. The supplied English TeX is a local translation/transcription,
not an original author TeX source. It has been read at 1.4.3--1.5.2 and
3.2.1--3.2.13; its compressed presentation of \S3.2 does not constitute
verification of all intermediate assertions in Deligne's original proof.
Every finite-dimensional assertion made below is proved here.

\subsection{The exact prime-power spectrum of the original source}

Retain the original integer $k\ge17$, $k\equiv1\pmod4$, $q=(k+1)^2$,
the original simple quartet with $\delta,\gamma>0$, the admitted period,
and each literal cutoff $q-1\le N\le2q$. Set
\begin{equation}\tag{TDP1}
\begin{gathered}
\lambda_{ab}=(2b-k)\gamma-i(2a-k)\delta,\qquad
Q_k(y)=\prod_{a,b=0}^k(y-\lambda_{ab}),\qquad
E=\mathbb C[y]/Q_k,\qquad M[f]=[yf],\\
s=k/2+iy,\qquad
F_p=\exp(i(\log p)M),\qquad
\widetilde F_p=p^{k/2}F_p .
\end{gathered}
\end{equation}
Here $p$ is a rational prime. The map $s=k/2+iy$ is the original
centred Mellin coordinate map; the scalar $p^{k/2}$ is retained in the
uncentred action. The parameters $k$ and $q$ have not been replaced by
the number of factors in any later tensor construction.

The $q$ roots in TDP1 are distinct: equality of two roots forces equality
of their real and imaginary parts, and hence $a=a'$ and $b=b'$.
The exact cardinal polynomials
\begin{equation}\tag{TDP2}
\ell_{ab}(y)=\prod_{(c,d)\ne(a,b)}
 \frac{y-\lambda_{cd}}{\lambda_{ab}-\lambda_{cd}}
\end{equation}
therefore form a basis of $E$. Evaluation at the roots proves
$M\ell_{ab}=\lambda_{ab}\ell_{ab}$ in $E$. Consequently, for every
integer $m\ge0$,
\begin{equation}\tag{TDP3}
\begin{aligned}
F_p^m\ell_{ab}
 &=p^{m(2a-k)\delta}
   \exp\!\bigl(im(2b-k)\gamma\log p\bigr)\ell_{ab},\\
\widetilde F_p^{\,m}\ell_{ab}
 &=p^{mk/2+m(2a-k)\delta}
   \exp\!\bigl(im(2b-k)\gamma\log p\bigr)\ell_{ab}.
\end{aligned}
\end{equation}
This follows by applying the convergent exponential series to each
cardinal vector. Different prime-action eigenvalues may coincide because
of their phases; the labelled vectors and their multiplicities remain.
In particular the exact largest eigenvalue modulus is $p^{mk\delta}$
for $F_p^m$ and $p^{mk/2+mk\delta}$ for $\widetilde F_p^{\,m}$.
These are statements about the programme's stipulated source, not
identifications of a finite trace with the scalar prime evaluations in
Weil's explicit formula.

Let $G_N$ be the complete original source metric, let
$K=\ker\Lambda$, and retain the original section
\begin{equation}\tag{TDP4}
\begin{gathered}
Q_B=(\Lambda G_N^{-1}\Lambda^*)^{-1},\qquad
L=G_N^{-1}\Lambda^*Q_B,\qquad J=LQ_B^{-1/2},\\
J^\dagger J=I_B,\qquad P=JJ^\dagger=L\Lambda,\qquad
Q=I_E-P,\qquad \Phi(X)=J^\dagger XJ.
\end{gathered}
\end{equation}
The dagger uses $G_N$ on $E$, its restriction on $K$, and the stated
Euclidean coordinates on $B$. Let $\iota_K:K\hookrightarrow E$ be the
isometric inclusion. The map $W:B\oplus K\to E$,
$W(y,z)=Jy+\iota_Kz$, is unitary for these metrics, since the two
summands are orthogonal, $\operatorname{im}J=K^\perp$, and their sum is
$E$. Define the full blocks
\begin{equation}\tag{TDP5}
W^\dagger MW=
\begin{pmatrix}H&B_{01}\\ B_{10}&D_K\end{pmatrix},\qquad
\begin{array}{ll}
H=J^\dagger MJ,&B_{01}=J^\dagger M\iota_K,\\
B_{10}=\iota_K^\dagger MJ,&D_K=\iota_K^\dagger M\iota_K.
\end{array}
\end{equation}
No invariance of $K$ under $M$ is assumed or used.

\subsection{The entire compression and the complete return through the kernel}

For every $z\in\mathbb C$ put $U(z)=J^\dagger e^{zM}J$ and
$Y(z)=\iota_K^\dagger e^{zM}J$. Differentiating the entire exponential
series and applying TDP5 gives
\begin{equation}\tag{TDP6}
\begin{gathered}
U'=HU+B_{01}Y,\qquad Y'=B_{10}U+D_KY,\qquad U(0)=I_B,\quad Y(0)=0,\\
Y(z)=\int_0^z e^{(z-t)D_K}B_{10}U(t)\,dt,\\
U'(z)=HU(z)+\int_0^z\mathcal K(z-t)U(t)\,dt,
\qquad \mathcal K(t)=B_{01}e^{tD_K}B_{10}.
\end{gathered}
\end{equation}
Integrals are along the straight segment, with the usual complex
orientation. The second equation follows by differentiating
$e^{-zD_K}Y(z)$ and integrating; substitution proves the third.
Thus TDP6 is an identity of entire functions, including purely imaginary
prime arguments. A second variation-of-constants calculation gives
\begin{equation}\tag{TDP7}
U(z)-e^{zH}=
\int_0^z\!\int_0^t
 e^{(z-t)H}B_{01}e^{(t-v)D_K}B_{10}U(v)\,dv\,dt.
\end{equation}
The difference is the full original-kernel return, including its complex
signs. No positive sign for this operator is asserted.

Insertion of $I_E=P+Q$ in $e^{(z+w)M}$ proves the exact product defect
\begin{equation}\tag{TDP8}
\begin{aligned}
U(z+w)-U(z)U(w)
 &=J^\dagger e^{zM}Qe^{wM}J,\\
U(mz)-U(z)^m
 &=\sum_{j=1}^{m-1}U(z)^{m-1-j}
   J^\dagger e^{zM}Qe^{jzM}J\qquad(m\ge1).
\end{aligned}
\end{equation}
For the second identity the sum is empty at $m=1$. If the assertion is
known for $m$, write
$U((m+1)z)-U(z)^{m+1}=U(z)(U(mz)-U(z)^m)
+U((m+1)z)-U(z)U(mz)$ and use the first identity. This proves it for
all integers $m$. Specializing $z=i\log p$ gives the exact correction
between the observed prime-power action and powers of the observed
one-prime action. For the uncentred action both sides acquire precisely
$p^{mk/2}$, with no change to any of the kernel factors.

The same defect determines the local determinant. Write
$W^\dagger F_pW=\left(\begin{smallmatrix}A_p&B_p\\C_p&D_p\end{smallmatrix}\right)$.
Over formal power series in $u$, every matrix with constant term $I$ is
invertible. Block Gaussian elimination then gives
\begin{equation}\tag{TDP9}
\begin{gathered}
\sum_{m\ge0}\Phi(F_p^m)u^m
 =\bigl[I_B-uA_p-u^2B_p(I_K-uD_p)^{-1}C_p\bigr]^{-1},\\
\det(I_E-uF_p)
 =\det(I_K-uD_p)\det(I_B-uA_p)\,\mathcal E_p(u),\\
\mathcal E_p(u)=
\det\!\left[I_B-u^2(I_B-uA_p)^{-1}
 B_p(I_K-uD_p)^{-1}C_p\right].
\end{gathered}
\end{equation}
To verify the first identity, expand $(I_E-uF_p)^{-1}$ as its geometric
series and take its upper-left block. Eliminate the lower-left block
using the invertible lower-right block $I_K-uD_p$; its Schur complement
is the displayed inverse. Taking determinants proves the second
identity, and factoring $I_B-uA_p$ proves the third. Since both sides are
rational functions, the identities also hold wherever the displayed
inverses exist, and extend across removable singularities after clearing
denominators. They retain every pole cancellation carried by $K$.

The formal logarithmic derivative, or the identity
$-\log\det(I-uT)=\sum_{m\ge1}\operatorname{tr}(T^m)u^m/m$, yields
\begin{equation}\tag{TDP10}
\begin{aligned}
-\log\mathcal E_p(u)
 &=\sum_{m\ge1}
 \frac{\operatorname{tr}(F_p^m)-\operatorname{tr}(A_p^m)
 -\operatorname{tr}(D_p^m)}{m}u^m,\\
[u^2](-\log\mathcal E_p(u))&=\operatorname{tr}(B_pC_p).
\end{aligned}
\end{equation}
Indeed the diagonal blocks of $F_p^2$ have traces
$\operatorname{tr}(A_p^2)+\operatorname{tr}(B_pC_p)$ and
$\operatorname{tr}(D_p^2)+\operatorname{tr}(C_pB_p)$; cyclicity of the
finite trace proves the coefficient claim. These traces are generally
complex. Positivity of $\Phi$ on positive operators does not change this
calculation into a nonnegative Euler coefficient.

\subsection{The support algebra exhibits the cancellation explicitly}

The original support representation and its boundary operator are
\begin{equation}\tag{TDP11}
\begin{gathered}
\rho(c_0,c_1+c_2\eta)=c_0(I-\Pi)+c_1\Pi+c_2R,\qquad
\Pi=ee^\dagger+ff^\dagger,\qquad R=\epsilon fe^\dagger,\\
e^\dagger f=0,\quad \|e\|=\|f\|=1,\quad \epsilon>0,\quad R^2=0,\\
[\tau]\longmapsto I_E,\qquad [n]\longmapsto\Pi+nR.
\end{gathered}
\end{equation}
The symbol multiplication in the last line represents addition in the
original monoid $\mathbb Z\sqcup\{\tau\}$. In particular $[p]^m=[mp]$
there, not $[p^m]$. The original $\tau$ action cannot be identified with
the multiplicative prime-power law by replacing one of these indices.
The exact one-parameter group in the same algebra is
$\rho(1,1+z\eta)=I_E+zR=e^{zR}$; at integral $z=n$ this is the image
of $[\tau]+[n]-[0]$. Its multiplication law follows from $R^2=0$.

Retain the measured columns and all their phases:
\begin{equation}\tag{TDP12}
\begin{gathered}
u=J^\dagger e,\quad v=J^\dagger f,\qquad
\begin{pmatrix}a&r\\\bar r&d\end{pmatrix}
 =[u,v]^*[u,v],\quad \alpha=1-a,\quad \mathfrak d=ad-|r|^2,\\
T=\Phi(R)=\epsilon vu^*,\qquad \kappa=\epsilon r,\qquad
T^2=\kappa T,\\
e_K=\iota_K^\dagger e,\quad f_K=\iota_K^\dagger f,\quad
\|e_K\|^2=\alpha,\quad\|f_K\|^2=1-d,\quad e_K^*f_K=-r.
\end{gathered}
\end{equation}
The identity $T^2=\kappa T$ is direct multiplication using $u^*v=r$.
The last line follows by inserting $P+Q$ in the inner products of $e,f$.
The support map preserves the original nilpotent on $E$, but compression
can give it the nonzero eigenvalue $\kappa$. Since $\mathfrak d>0$ in
the original finite area guards, neither $u$ nor $v$ is zero, so $T\ne0$.

The complete $R$ blocks, its memory and its two exponentials are
\begin{equation}\tag{TDP13}
\begin{gathered}
W^\dagger RW=
\begin{pmatrix}T&\epsilon ve_K^*\\
 \epsilon f_Ku^*&\epsilon f_Ke_K^*\end{pmatrix},\qquad
\mathcal K_R(z)=-\kappa e^{-\kappa z}T,\\
\Phi(e^{zR})=I_B+zT,\qquad
e^{zT}=I_B+f_\kappa(z)T,\qquad
f_\kappa(z)=\begin{cases}(e^{\kappa z}-1)/\kappa,&\kappa\ne0,\\
z,&\kappa=0.
\end{cases}
\end{gathered}
\end{equation}
For the memory identity, the lower-right block sends $f_K$ to
$-\kappa f_K$. Thus its exponential sends $f_K$ to
$e^{-\kappa z}f_K$, and multiplication by the two off-diagonal blocks
gives $-\kappa e^{-\kappa z}T$. The first exponential truncates because
$R^2=0$. For the second, induction gives $T^j=\kappa^{j-1}T$ for $j\ge1$;
substitution in the entire exponential series proves the formula,
including its removable value at $\kappa=0$.

For direct verification that the memory cancels the extra exponential,
substitute $U_R(z)=I_B+zT$ in TDP6 with the blocks in TDP13. The sum of
the two terms on its right is
$T+z\kappa T-\kappa T\int_0^z e^{-\kappa(z-t)}(1+t\kappa)dt=T$,
since the integrand is the derivative in $t$ of
$t e^{-\kappa(z-t)}$. This is precisely $U_R'(z)$.

Put $c=i\log p$, $n_E=\dim E=q$ and $b_B=\dim B$. The exact support
prime action $I_E+cR$ and all its prime powers satisfy
\begin{equation}\tag{TDP14}
\begin{gathered}
\Phi((I_E+cR)^m)=I_B+mcT,\\
(I_B+cT)^m=I_B+
\begin{cases}((1+c\kappa)^m-1)T/\kappa,&\kappa\ne0,\\
mcT,&\kappa=0,
\end{cases}\\
\det(I_E-w(I_E+cR))=(1-w)^{n_E}.
\end{gathered}
\end{equation}
The first line is the binomial theorem with $R^2=0$. For the second line
expand by the binomial theorem and use $T^j=\kappa^{j-1}T$; the scalar
binomial formula supplies the coefficient. All eigenvalues of $R$ are
zero, so all eigenvalues of $I_E+cR$ equal one; this proves the last line.
In contrast, the observed block has a candidate eigenvalue
$1+c\kappa$. This is a compression effect, not a moved eigenvalue of
the full support action.

When $K\ne0$, its original block is $I_K+c\epsilon f_Ke_K^*$ and has
the candidate eigenvalue $1-c\kappa$. Its exact factors in TDP9 are
\begin{equation}\tag{TDP15}
\begin{aligned}
\det(I_B-w(I_B+cT))
 &=(1-w)^{b_B-1}(1-w-wc\kappa),\\
\det(I_K-w(I_K+c\epsilon f_Ke_K^*))
 &=(1-w)^{n_E-b_B-1}(1-w+wc\kappa),\\
\mathcal E_{p,R}(w)
 &=\frac{(1-w)^2}{(1-w)^2-w^2c^2\kappa^2}.
\end{aligned}
\end{equation}
These identities first hold formally at $w=0$. The first two follow
from the rank-one determinant identity
$\det(I+xy^*)=1+y^*x$, which is proved by extending $x$ to a basis when
$x\ne0$; the case $x=0$ is immediate. Their product, the last line and
TDP14 multiply to an identity. Alternatively substituting the blocks
of TDP13 into TDP9 gives the last line directly. In the inverse
determinant, its factor $\mathcal E_{p,R}(w)^{-1}$ cancels the two
candidate poles at $1/(1+c\kappa)$ and $1/(1-c\kappa)$. If $\kappa=0$
this factor is identically one. If $K=0$, then $r=0$, $\kappa=0$ and the
kernel factor is absent; TDP14 remains valid. This is a complete
pole-cancellation calculation on the actual support receiver.

The original estimates do not authorize deleting $\kappa$:
\begin{equation}\tag{TDP16}
\begin{gathered}
|\kappa|^2=\epsilon^2|r|^2
 \le\epsilon^2\alpha(1-d),\qquad
\|T\|=\epsilon\sqrt{ad},\\
\alpha\le e^{-2q\psi(t_N)+O_{h,\varpi}(k\log(q+2))},\quad
\log\epsilon=q\psi(t_N)+O_{h,\varpi}(k\log(q+2)),\\
|\kappa|\le e^{O_{h,\varpi}(k\log(q+2))}\sqrt{1-d},
\qquad t_N=(N+1-q)/q.
\end{gathered}
\end{equation}
The first inequality is the Cauchy--Schwarz inequality on $K$ using
TDP12; the rank-one norm gives the second equality. Substitution of the
two original estimates proves the last line. The large rate in
$\epsilon$ cancels the square-root leakage rate, and the phase of $r$
has not been restricted. Therefore the displayed estimate supplies
neither a sign nor a vanishing assertion for $\kappa$.

The full $M=C+R$ in ST3 must still be used. For $\lambda$ outside the
spectrum of the original self-adjoint $C$, rank-one elimination gives
\begin{equation}\tag{TDP17}
\det(\lambda I_E-M)=\det(\lambda I_E-C)
 \bigl[1-\epsilon e^\dagger(\lambda I_E-C)^{-1}f\bigr].
\end{equation}
Factor $\lambda I-C$ on the left and use the rank-one determinant
identity just proved. Clearing denominators extends the equality to a
polynomial identity. Thus nilpotence of $R$ by itself does not remove
the spectral displacement in TDP3. In particular products
$RC^jR=\epsilon^2(e^\dagger C^jf)fe^\dagger$ need not vanish.
The exact interaction with $C$, not $R$ in isolation, determines the
original source spectrum.

\subsection{Tensor products retain the same defect, with exact mixed sectors}

Let $\nu\ge1$ be an independent tensor count and retain the product
metric $G_N^{\otimes\nu}$. Define
\begin{equation}\tag{TDP18}
\begin{gathered}
J_\nu=J^{\otimes\nu},\quad P_\nu=P^{\otimes\nu},\quad
Q_\nu=I-P_\nu
 =\sum_{\varnothing\ne S\subset\{1,\ldots,\nu\}}
 \bigotimes_{j=1}^{\nu}\begin{cases}Q,&j\in S,\\P,&j\notin S,
 \end{cases}\\
M_\nu=\sum_{j=1}^\nu I^{\otimes(j-1)}\otimes M\otimes
 I^{\otimes(\nu-j)},\qquad H_\nu=\sum_{j=1}^\nu
 I^{\otimes(j-1)}\otimes H\otimes I^{\otimes(\nu-j)},\\
J_\nu^\dagger e^{zM_\nu}J_\nu=U(z)^{\otimes\nu},\qquad
e^{zH_\nu}=(e^{zH})^{\otimes\nu}.
\end{gathered}
\end{equation}
The projections in the first line have mutually orthogonal ranges:
two distinct subsets place $PQ=0$ in at least one factor. Expansion of
$(P+Q)^{\otimes\nu}$ proves their sum. The summands of $M_\nu$ commute,
so multiplication of their convergent exponential series proves the
last line. Tensoring therefore preserves the full source action, and
also preserves the difference caused by observation:
\begin{equation}\tag{TDP19}
U(z)^{\otimes\nu}-(e^{zH})^{\otimes\nu}
 =\sum_{j=1}^\nu U(z)^{\otimes(j-1)}\otimes
 \bigl(U(z)-e^{zH}\bigr)\otimes(e^{zH})^{\otimes(\nu-j)}.
\end{equation}
The right side telescopes after inserting and subtracting tensors with
successively one further factor $U(z)$. No mixed-kernel subset is
discarded by TDP18 or TDP19.

For the original support put $A=\Phi(\Pi)=[u,v][u,v]^*$, with its two
strictly positive eigenvalues $\vartheta_+,\vartheta_-$ from ST7. Let
$A_\nu=A^{\otimes\nu}$, which is the exact compression of
$\Pi^{\otimes\nu}$. Its eigenvalues and defect are
\begin{equation}\tag{TDP20}
\begin{gathered}
\vartheta_+^{\nu-j}\vartheta_-^j
 \quad\hbox{with multiplicity }{\nu\choose j},\quad 0\le j\le\nu,
 \quad\hbox{and }0\hbox{ on the remaining }b_B^\nu-2^\nu\hbox{ dimensions},\\
A_\nu-A_\nu^2
 =\sum_{j=1}^{\nu}(A^2)^{\otimes(j-1)}\otimes(A-A^2)
 \otimes A^{\otimes(\nu-j)}
 =(Q_\nu\Pi^{\otimes\nu}J_\nu)^\dagger
  (Q_\nu\Pi^{\otimes\nu}J_\nu).
\end{gathered}
\end{equation}
Tensor products of orthonormal eigenvectors prove the first line,
including labelled multiplicities when values coincide. A telescoping
expansion proves the first equality in the second line. Since
$\Pi^{\otimes\nu}$ is an orthogonal projection, insertion of
$P_\nu+Q_\nu$ proves the last equality. Each summand of the middle
expression is positive semidefinite. Positivity of this defect does
not make it zero.

There are exact amplification values, not merely an unspecified error:
\begin{equation}\tag{TDP21}
\begin{gathered}
\det{}^+A_\nu=\mathfrak d^{\,\nu2^{\nu-1}},\qquad
\frac{\log\det{}^+A_\nu}{\nu2^\nu}=\tfrac12\log\mathfrak d,\\
\left|\frac{\log\det{}^+A_\nu}{\nu2^\nu}
 -\tfrac12\log d\right|\le\frac{\alpha}{2\Theta_N},\qquad
0\le1-\vartheta_+^\nu\le\nu\alpha,\\
\frac1\nu\log(\vartheta_+^{\nu-j}\vartheta_-^j)
 =(1-j/\nu)\log\vartheta_++(j/\nu)\log\vartheta_-.
\end{gathered}
\end{equation}
Each of $\vartheta_+$ and $\vartheta_-$ occurs $\nu2^{\nu-1}$ times in
the product of the $2^\nu$ nonzero eigenvalues, proving the determinant
identity. ST11 then proves its estimate. The identity
$1-x^\nu=(1-x)\sum_{j=0}^{\nu-1}x^j$ for $0\le x\le1$, together with
ST8, proves the bound on the first tensor branch. The final line is the
logarithm of the exact eigenvalue. In particular a non-unit second
support eigenvalue contributes a loss proportional to the tensor
count. Taking a $\nu$th root does not remove that value. This is the
precise contrast with a fixed exponent error divided by the tensor
count in Deligne's amplification.

The support representation survives observation as a homomorphism
exactly when both original support vectors are retained:
\begin{equation}\tag{TDP22}
\Phi\circ\rho\text{ is multiplicative}\quad\Longleftrightarrow\quad
a=d=1\quad\Longleftrightarrow\quad e,f\in K^\perp.
\end{equation}
For necessity, multiplicativity makes $A^2=A$. Its two nonzero
eigenvalues must then be one. Thus $a+d=\operatorname{tr}A=2$ and
$0\le a,d\le1$ imply $a=d=1$; Cauchy--Schwarz on $K$ gives $r=0$.
The norm identities in TDP12 imply $e,f\in K^\perp$. Conversely, in
that case both $\Pi$ and $R$ preserve $K^\perp$ and are zero on $K$.
Their compressions obey the identical multiplication table
$A^2=A$, $AT=TA=T$, $T^2=0$. The complement $I_B-A$ has zero products
with $A,T$, so the algebra formula in TDP11 is multiplicative.
The defect space determined by failure is the actual range
$\operatorname{im}(Q\Pi J)\subset K$, with Gram operator $A-A^2$,
and its tensor analogues are the explicitly decomposed ranges in TDP20.

The support nilpotent also tensors without dropping mixed terms. Put
$R_\nu=\sum_{j=1}^\nu I^{\otimes(j-1)}\otimes R\otimes I^{\otimes(\nu-j)}$.
Commutativity between different factors and $R^2=0$ give
\begin{equation}\tag{TDP23}
\begin{gathered}
R_\nu^j=j!\sum_{|S|=j}\bigotimes_{l=1}^\nu
 \begin{cases}R,&l\in S,\\I_E,&l\notin S,\end{cases}
\quad (0\le j\le\nu),\qquad
R_\nu^{\nu+1}=0,\quad R_\nu^\nu=\nu!R^{\otimes\nu}\ne0,\\
e^{zR_\nu}=(I_E+zR)^{\otimes\nu},\qquad
\Phi_\nu(e^{zR_\nu})=(I_B+zT)^{\otimes\nu},\qquad
\Phi_\nu(X)=J_\nu^\dagger XJ_\nu.
\end{gathered}
\end{equation}
In expanding the $j$th power, only words using $j$ distinct positions
survive; every subset occurs in $j!$ orders. Since $R\ne0$, choose $x$
with $Rx\ne0$; then $R^{\otimes\nu}x^{\otimes\nu}\ne0$, proving its
nonvanishing. Thus the full support tensor has only eigenvalue one
after exponentiation, while its nilpotent length is exactly $\nu+1$.
Observation keeps the explicitly displayed polynomial in $z$ of degree
at most $\nu$. It cannot be replaced by $e^{z\Phi_\nu(R_\nu)}$ without
the memory already calculated.

\subsection{The exact object carrying eigenvalues through observation}

For each $\lambda\in\mathbb C$, define the lifted observed eigenspace
and its observed image by
\begin{equation}\tag{TDP24}
\begin{gathered}
\mathscr L_\lambda=
 \{(y,z)\in B\oplus K:
 (H-\lambda I)y+B_{01}z=0,
 \ B_{10}y+(D_K-\lambda I)z=0\},\\
\mathscr O_\lambda=\{y\in B:\text{there is }z\in K
 \text{ with }(y,z)\in\mathscr L_\lambda\},\\
0\longrightarrow\ker(D_K-\lambda I)\cap\ker B_{01}
 \xrightarrow{z\mapsto(0,z)}\mathscr L_\lambda
 \xrightarrow{(y,z)\mapsto y}\mathscr O_\lambda\longrightarrow0,\\
\mathscr L_\lambda\xrightarrow[\cong]{\ W\ }\ker(M-\lambda I),
\qquad\|W(y,z)\|_{G_N}^2=\|y\|^2+\|z\|^2.
\end{gathered}
\end{equation}
Multiplication by the block matrix in TDP5 proves the last isomorphism.
The first map in the exact sequence is injective; the kernel of the
second consists exactly of pairs $(0,z)$ satisfying both displayed
equations; and surjectivity is the definition of $\mathscr O_\lambda$.
Orthogonality in TDP4 proves the metric identity. If
$\lambda\notin\operatorname{spec}D_K$, this same space is the explicit
graph
\begin{equation}\tag{TDP25}
\begin{gathered}
z=(\lambda I-D_K)^{-1}B_{10}y,\qquad
\mathscr O_\lambda=
\ker\bigl[H+B_{01}(\lambda I-D_K)^{-1}B_{10}-\lambda I_B\bigr],\\
\|(y,z)\|^2=
y^*\bigl[I_B+B_{10}^\dagger(\bar\lambda I-D_K^\dagger)^{-1}
 (\lambda I-D_K)^{-1}B_{10}\bigr]y.
\end{gathered}
\end{equation}
Solving the second equation of TDP24 proves the first line, and taking
the squared norm proves the second. At exceptional $\lambda$ TDP24
remains the exact space, with no discarded singular fibre.

For a nonzero original eigenvector $x$, set $y=J^\dagger x$ and
$z=\iota_K^\dagger x$. Even when the programme proves $y\ne0$ for its
terminal class, the exact observed equation is
$(H-\lambda I)y=-B_{01}z$, not an eigenvalue equation for $H$ alone.
Tensor products preserve the lift: $x^{\otimes\nu}$ maps to
$y^{\otimes\nu}$ in the fully observed sector, and to every other
subset sector of TDP18 with the appropriate $y$ or $z$ factor.
If $y\ne0$, its exact detection ratio is
\begin{equation}\tag{TDP26}
\frac{\|J_\nu^\dagger x^{\otimes\nu}\|}{\|x^{\otimes\nu}\|}
 =\left(\frac{\|y\|}{\sqrt{\|y\|^2+\|z\|^2}}\right)^\nu.
\end{equation}
This follows from the definition of the Hilbert tensor metric and
TDP24. A fixed non-unit detection factor thus gives a linear logarithmic
cost in the tensor count; it is not a fixed additive error that disappears
after taking a tensor root. If $y=0$, the full lift still exists but its
fully observed tensor sector is zero. No lower bound for that case has
been invented.

\subsection{What this changes in a Deligne-style amplification}

Deligne's two relevant mechanisms have different exact inputs. In
1.4.3--1.5.2 the cohomological denominator has a bounded weight shift;
the trace formula identifies it with the Euler product; nonnegative
coefficients for the even tensor powers prevent a local pole from
being hidden inside the product. The finite tensor calculation
$w(\alpha^{2\nu})=2\nu w(\alpha)$ then divides the fixed error by
$2\nu$. In the supplied version of 3.2, an eigenvalue on $H^1_c$
squares inside a K\"unneth component of a surface, while the surface
estimate improves the induction error. Its displayed inequalities have
the arithmetic implication
\begin{equation}\tag{TDP27}
2w(\alpha)<2+2^{-j}\quad\Longrightarrow\quad
w(\alpha)<1+2^{-j-1}.
\end{equation}
This line is just division by two. Attribution of the full surrounding
proof remains to Deligne; the local abbreviated English source is not
used to supply omitted hypotheses or maps.

For the present programme, TDP3 gives the actual prime eigenvalues,
TDP24 gives their exact observation map, TDP18 gives the full tensor
object, and TDP9 gives the exact local determinant after observation.
The kernel contribution is the concrete rational function
$\mathcal E_p$, not an unnamed missing term. Its support-only value in
TDP15 cancels precisely the false poles that discarding the kernel
would introduce. Accordingly the positivity of $A-A^2$ in ST9 cannot
replace the noncancellation input in Deligne's pole argument: TDP15 is
an exact example with a positive support defect and nontrivial pole
cancellation. Likewise TDP21 and TDP26 calculate the tensor loss itself;
the second support eigenvalue and the eigenvector's unobserved component
are retained as spaces and maps, not removed by a weight declaration.

All formulas apply separately at $N=q-1,q,2q-1,2q$. Any signed return
uses the original signs $+,+,-,-$; for scalar absolute error estimates
the triangle inequality gives the sum of the four respective errors.
No cancellation between these cutoffs has been assumed. No equivalence
between the original conductor and Deligne's full cohomological trace
formula, and no conclusion about RH, is claimed from a compressed
positivity assertion. The calculation instead supplies the exact
prime-power, determinant and tensor defect objects that the original
programme now has to evaluate, with their connecting maps proved.

There is an exact positive receiver for every finite original prime
combination, with the adjoint in the original metric. Let $\mathcal I$
be any finite set of prime-power indices $(p,m)$, let $c_{p,m}\in\mathbb C$
be arbitrary, and put $X=\sum_{(p,m)\in\mathcal I}c_{p,m}F_p^m$. Then
\begin{equation}\tag{TDP28}
\begin{aligned}
\Phi(X^\dagger X)-\Phi(X)^\dagger\Phi(X)
 &=(QXJ)^\dagger(QXJ)\\
 &=\sum_{(p,m),(r,n)\in\mathcal I}
 \overline{c_{p,m}}c_{r,n}\,
 J^\dagger(F_p^m)^\dagger QF_r^nJ\succeq0.
\end{aligned}
\end{equation}
Insertion of $I_E=P+Q$ between $X^\dagger$ and $X$ proves the first
identity, since $Q^2=Q=Q^\dagger$ and $P=JJ^\dagger$. Expansion proves
the second identity, with all phases retained. Evaluation on $b\in B$
gives $\|QXJb\|_{G_N}^2$, proving positivity. This is an exact
prime-action Gram identity. Equating it to Weil's scalar positivity
functional would require an additional proved map; no such equality is
asserted here. The positive adjoint-product defect and the Euler
product defect in TDP9 are both calculated, and they have different
factors: the latter contains $B_p(I-uD_p)^{-1}C_p$, without an adjoint.
